\section{Principios Éticos Aplicados}

Cuatro principios estructuran la evaluación ética de Astro-Denoise. Cada uno aborda un riesgo distinto y se materializa en decisiones técnicas concretas.

\subsection{Reproducibilidad}

La reproducibilidad es la capacidad de obtener los mismos resultados con los mismos datos y métodos. En astronomía es indispensable: las observaciones son únicas en el tiempo y las decisiones científicas implican inversiones enormes.

Astro-Denoise la garantiza mediante código abierto, datos versionados (LSST DC2), configuración explícita, semillas registradas y resultados en formato auditable. Cualquier investigador puede clonar el repositorio y verificar los números. Esto no es aspiracional, es garantía técnica (\cite{harrison2017}).

\subsection{Equidad}

La equidad significa igualdad de oportunidad, no de resultados. La ciencia astronómica de punta hoy está concentrada en instituciones del Norte Global con acceso a consorcios como LSST o Euclid. Investigadores en América Latina enfrentan barreras técnicas, idiomáticas y económicas.

Astro-Denoise las elimina: funciona en laptops, documenta en español, requiere aprobación de nadie, y valida contribuciones con el mismo rigor para todos. El impacto esperado es que investigadores del Sur Global publiquen, innoven y validen ciencia local sin depender de consorcios internacionales.

\subsection{Gobernanza de Datos}

Gobernanza significa claridad sobre responsabilidades, trazabilidad de origen y cumplimiento normativo. Astro-Denoise opera con licencias explícitas (MIT para código, CC0 para datos, CC-BY para documentación), metadatos completos de la simulación LSST DC2, y un plan de preservación digital. La Universidad La Salle asume responsabilidad institucional.

\subsection{Excelencia Científica}

Excelencia es aplicar estándares rigurosos sin importar la procedencia institucional. En Astro-Denoise se manifiesta en métricas múltiples e independientes, revisión por pares del código, y transparencia total en el proceso de validación. No hay ventajas injustas: el mismo estándar para todos.

\subsection{Integración}

Los cuatro principios forman un sistema: reproducibilidad es el \textbf{método}, equidad es el \textbf{objetivo}, gobernanza es el \textbf{mecanismo}, y excelencia es el \textbf{estándar}. Juntos, crean las condiciones para que cualquier investigador, en cualquier país, pueda confiar en los resultados del proyecto.
